\subsection{应用：离散赋值环的新刻画}

命题 11. 设 A 是局部 Krull 整环（特别地，是整闭的局部诺特整环），m 是其极大理想。以下条件等价：

\begin{enumerate}
\def\labelenumi{(\alph{enumi})}
\item
  A 是离散赋值环；
\item
  m 是可逆的；
\item
  A: m ≠ A；
\item
  m 是除性的；
\item
  m 是 A 的唯一非零素理想。
\end{enumerate}

由于离散赋值环的每个非零理想都是主理想（第 VI 章，§ 3，第 6 节，命题 9），它是可逆的，从而 (a) 蕴含 (b)。若 m 可逆，其逆是 A: m（第 II 章，§ 5，第 6 节，命题 10），因此 A: m ≠ A；从而 (b) 蕴含 (c)。若 A: m ≠ A，则 A: (A: m) ≠ A；现在 m ⊂ A: (A: m)；由于 m 是极大理想，故 m = A: (A: m)，所以 m 是除性的（第 1 节，命题 1 (c)）；从而 (c) 蕴含 (d)。(d) 蕴含 (e) 由第 6 节的定理 3 得出。最后，若 m 是 A 的唯一非零素理想，则其高度为 1，因而 A \(_{m}\) 是离散赋值环（第 6 节，定理 4）；由于 A 是局部环，A \(_{m}\) = A，这说明 (e) 蕴含 (a)。

